\section{Lazy Segment Tree}
\whenuse{Actualizaciones y consultas sobre intervalos, por ejemplo sumar en un
rango y consultar suma/m\'inimo/m\'aximo sin tocar cada elemento.}

Suma en rango y suma sobre rango. Los intervalos son semiabiertos $[l,r)$.
Para consultar m\'inimos o m\'aximos manteniendo las sumas en rango, reemplazar
\texttt{NEUTRAL}, \texttt{combine} y la l\'inea marcada en \texttt{apply}
por el mismo caso comentado. No cambiar \texttt{push}, build, add ni query.
Cada nodo combina sus hijos; \texttt{lazy[p]} guarda el incremento pendiente
por elemento. El m\'inimo/m\'aximo aumenta en $v$, pero la suma en $v(r-l)$.

\complexity{build $O(n)$; update/query $O(\log n)$; memoria $O(n)$.}

\begin{lstlisting}[style=cpp]
struct LazySegmentTree {
    // SUMA (activo). MIN/MAX: cambiar estas dos lineas y apply.
    static constexpr ll NEUTRAL = 0;
    static ll combine(ll a, ll b) { return a + b; }
    // MIN: static constexpr ll NEUTRAL = numeric_limits<ll>::max();
    // MIN: static ll combine(ll a, ll b) { return min(a, b); }
    // MAX: static constexpr ll NEUTRAL = numeric_limits<ll>::min();
    // MAX: static ll combine(ll a, ll b) { return max(a, b); }

    int n;
    vector<ll> tree, lazy;

    explicit LazySegmentTree(const vector<ll> &a)
        : n(a.size()), tree(4 * max(1, n)), lazy(4 * max(1, n)) {
        if (n) build(1, 0, n, a);
    }

    void build(int p, int l, int r, const vector<ll> &a) {
        if (r - l == 1) { tree[p] = a[l]; return; }
        int m = (l + r) / 2;
        build(p << 1, l, m, a);
        build(p << 1 | 1, m, r, a);
        tree[p] = combine(tree[p << 1], tree[p << 1 | 1]);
    }

    void apply(int p, int l, int r, ll value) {
        tree[p] += value * (r - l); // SUMA (activo).
        // tree[p] += value;       // MIN/MAX: reemplaza la anterior.
        lazy[p] += value; // Igual para SUMA, MIN y MAX.
    }

    void push(int p, int l, int r) {
        if (lazy[p] == 0 || r - l == 1) return;
        int m = (l + r) / 2;
        apply(p << 1, l, m, lazy[p]);
        apply(p << 1 | 1, m, r, lazy[p]);
        lazy[p] = 0;
    }

    void add(int ql, int qr, ll value) {
        if (n) add(1, 0, n, ql, qr, value);
    }
    void add(int p, int l, int r, int ql, int qr, ll value) {
        if (qr <= l || r <= ql) return;
        if (ql <= l && r <= qr) { apply(p, l, r, value); return; }
        push(p, l, r);
        int m = (l + r) / 2;
        add(p << 1, l, m, ql, qr, value);
        add(p << 1 | 1, m, r, ql, qr, value);
        tree[p] = combine(tree[p << 1], tree[p << 1 | 1]);
    }

    ll query(int ql, int qr) {
        return n ? query(1, 0, n, ql, qr) : NEUTRAL;
    }
    ll query(int p, int l, int r, int ql, int qr) {
        if (qr <= l || r <= ql) return NEUTRAL;
        if (ql <= l && r <= qr) return tree[p];
        push(p, l, r);
        int m = (l + r) / 2;
        return combine(query(p << 1, l, m, ql, qr),
                       query(p << 1 | 1, m, r, ql, qr));
    }
};
\end{lstlisting}

Con $n>0$, \texttt{tree[1]} contiene la respuesta global en $O(1)$.
\texttt{query(l,r)} cuesta $O(\log n)$ y devuelve el neutro si $l=r$.
Las sumas de valores y los acumulados lazy deben entrar en \texttt{ll}.

\warning{Estas variantes conservan la actualizaci\'on de sumar en rango.
No basta cambiar \texttt{combine} para consultar GCD, producto u operaciones
bit a bit: su resumen no permite aplicar una suma usando s\'olo este lazy.}

Para asignaci\'on de rango se necesita un lazy adicional que distinga ``sin
asignaci\'on'' de asignar cero.

\subsection{Reconocer el patr\'on: costos por umbral}
\begin{enumerate}
\item \textbf{Reordenar sin empeorar.} Si hay varias operaciones y se pide
minimizar su cantidad, preguntar si se pueden agrupar por tipo. Demostrar
que intercambiarlas conserva la validez y no aumenta el costo; no asumir
que siempre conmutan. Si funciona, contar cu\'antas veces usar cada tipo
puede reemplazar la simulaci\'on paso a paso.
\item \textbf{Fijar un umbral.} Si una operaci\'on afecta a todos los objetos
y otra a uno solo, fijar $x$ usos de la global y calcular el costo restante.
Buscar una funci\'on $f(x)=g(x)+\sum_i h_i(x)$, donde cada objeto aporta
seg\'un su relaci\'on con $x$. Probar que el costo calculado es necesario y
que existe una estrategia que lo alcanza.
\item \textbf{Acotar los candidatos.} Obtener una soluci\'on factible de
costo $B$. Si $f(x)\ge x$, ning\'un $x>B$ mejora esa soluci\'on. Esta
desigualdad, no el tama\~no de los valores de entrada, justifica recorrer
s\'olo $[0,B]$. Si no hay una cota peque\~na, buscar puntos donde cambia la
funci\'on y evaluar si se pueden comprimir coordenadas.
\item \textbf{Traducir cada cambio a un rango.} Si insertar un objeto suma
un mismo valor a $f(x)$ para $x<c$, actualiza el prefijo $[0,c)$; si afecta
a $x\ge c$, actualiza un sufijo. Borrar revierte esa contribuci\'on. Para
sumar sobre intervalos y pedir m\'inimos/m\'aximos, usar este lazy con el
resumen adecuado. Si s\'olo cambian posiciones individuales, alcanza el
Segment Tree puntual. Verificar siempre que el resumen permita aplicar
la actualizaci\'on sin recorrer las hojas.
\end{enumerate}

\subsection{Ejemplo: Emergency Rations}
Una caja tiene $c$ raciones. Cada d\'ia se vac\'ia una caja o se quita una
raci\'on de todas las no vac\'ias. Tras cada alta/baja, pedir el m\'inimo de
d\'ias, sin consumir realmente las raciones.

Con $x$ operaciones globales, cada caja con $c>x$ necesita una operaci\'on
individual. Hacer primero las globales y luego vaciar esas cajas alcanza
el costo (si todas se vac\'ian antes, se omiten los d\'ias sobrantes).
As\'i, el \textbf{m\'inimo} buscado es
\[
\min_{0\le x\le Q} f(x),\qquad f(x)=x+C(x),\quad
C(x)=\#\{\text{cajas con }c>x\}.
\]
Hay a lo sumo $Q$ cajas: vaciarlas individualmente cuesta a lo sumo $Q$,
y $f(x)\ge x$ permite descartar $x>Q$. Inicialmente $f(x)=x$.
Agregar una caja suma $1$ en $[0,\min(c,Q+1))$; quitarla resta $1$ all\'i.

\textbf{Antes de usar el ejemplo}, configurar el struct anterior para
\textbf{MIN}: \texttt{NEUTRAL = numeric\_limits<ll>::max()},
\texttt{combine(a,b) = min(a,b)} y, en \texttt{apply},
\texttt{tree[p] += value} sin multiplicar por la longitud.
Build y add deben combinar siempre con m\'inimo.
El namespace permite conservar el ejemplo junto al template; al enviar,
copiar solamente su \texttt{main} y el struct configurado para MIN.

\begin{lstlisting}[style=cpp]
namespace emergency_rations_example {
int main() {FIN;
    int Q;
    cin >> Q;

    vector<ll> a(Q + 1);
    iota(a.begin(), a.end(), 0LL);
    LazySegmentTree st(a); // Configurado para MIN.

    for (int i = 0; i < Q; ++i) {
        ll x;
        cin >> x;

        ll c = abs(x);
        int fin = min(c, (ll)Q + 1);
        // [0, fin): candidatos menores que c.
        st.add(0, fin, x > 0 ? 1 : -1);

        cout << st.tree[1]
             << (i + 1 == Q ? '\n' : ' ');
    }
    return 0;
}
}
\end{lstlisting}

Para \texttt{+1 +2 +2 +7 -2 +10 -1 -2}, las respuestas son
\texttt{1 2 2 3 3 4 3 2}. Cada cambio cuesta $O(\log Q)$ y leer la
ra\'iz cuesta $O(1)$; tiempo total $O(Q\log Q)$ y memoria $O(Q)$.
